\documentclass[11pt,letterpaper]{article}
\usepackage{technical-design-spec}
% Shared non-semantic helper definitions; one semantic module is loaded by each root.
\newcommand{\PTSource}[1]{\par\smallskip{\footnotesize\color{Meta}\textbf{Book crosswalk:} #1}\par\smallskip}
\newcommand{\PTOutcome}[1]{\begin{tipbox}[Competency outcome]#1\end{tipbox}}
\newcommand{\PTEvidence}[1]{\begin{notebox}[Required evidence]#1\end{notebox}}
\newcount\ptlinecount
\newcommand{\PTLines}[1]{\par\ptlinecount=0\loop\ifnum\ptlinecount<#1\par\noindent\rule{\linewidth}{0.25pt}\vspace{0.22cm}\advance\ptlinecount by1\repeat\par}
\newcommand{\PTField}[1]{\noindent\textbf{#1}\quad\rule{0.66\linewidth}{0.3pt}\par\vspace{0.22cm}}
\newcommand{\PTRating}{\noindent\textbf{Rating:}\quad 0\quad 1\quad 2\quad 3\quad 4\qquad\textbf{Evidence ref:}\rule{0.30\linewidth}{0.3pt}\par\vspace{0.14cm}}
\newcommand{\PTTask}[2]{\par\medskip\noindent\textbf{Task #1.} #2\par}

\setDocTitle{Chapter-to-Engineering Crosswalk}
\setDocSubtitle{All 18 source chapters, prerequisites, practical work and provenance}
\setDocVersion{1.0}\setDocStatus{Engineering education / source-aligned}
\setDocOwner{Compiler Engineering}
\setDocDate{30 October 2026}
\setDocSummary{Engineering reference and source crosswalk.}
\begin{document}\makedoctitle\makedocsummary
\tableofcontents\clearpage

\section{Interpretation}
Chapter numbers and printed-page starting points are navigational locators into Grune and Jacobs (2008). Original laboratory tasks and current security controls are not claims made by the textbook. The original book contains 18 chapters; Chapter 18 is an annotated bibliography.\par
\small\begin{longtable}{p{0.05\linewidth}p{0.32\linewidth}p{0.16\linewidth}p{0.34\linewidth}}\toprule Ch. & Chapter title & Prerequisite & Primary evidence\\\midrule\endhead
1 & Introduction & Recursion & Draw the dataflow from input text through tokens, concrete syntax tree, abstract syntax tree, and semantic checks.\\
2 & Grammars as a Generating Device & Recursion & Implement a reachability and productivity pass over nonterminals; report unused productions.\\
3 & Introduction to Parsing & Recursion & Draw two parse trees; design a shared parse-forest representation and acceptance criterion.\\
4 & General Non-Directional Parsing & Ch. 3 & Implement a triangular dynamic-programming table and preserve split backpointers.\\
5 & Regular Grammars and Finite-State Automata & Ch. 4 & Produce state transition tables and acceptance tests including empty and truncated prefixes.\\
6 & General Directional Top-Down Parsing & Ch. 5 & Implement a recursive-descent expression recognizer and demonstrate its behavior on a left-recursive grammar.\\
7 & General Directional Bottom-Up Parsing & Ch. 6 & Implement an Earley recognizer and compare its acceptance results with CYK for a compatible grammar.\\
8 & Deterministic Top-Down Parsing & Ch. 7 & Construct an LL(1) table with conflict detection and demonstrate a successful table trace.\\
9 & Deterministic Bottom-Up Parsing & Ch. 8 & Generate canonical item sets and record table conflicts for an ambiguous arithmetic grammar.\\
10 & Non-Canonical Parsers & Ch. 9 & Write a comparison report identifying which derivation restrictions are relaxed and what is recovered.\\
11 & Generalized Deterministic Parsers & Ch. 10 & Design branch-merging and shared-forest invariants; implement a bounded demonstration or compare an existing tool.\\
12 & Substring Parsing & Ch. 11 & Specify a substring query API and construct fixtures for full matches, prefixes, suffixes and true interior fragmen\\
13 & Parsing as Intersection & Ch. 12 & Build a product-style recognition exercise and record intermediate grammar size.\\
14 & Parallel Parsing & Ch. 13 & Design a parallel schedule and quantify work, depth, communication, and memory bandwidth.\\
15 & Non-Chomsky Grammars and Their Parsers & Ch. 14 & Create a comparison matrix for at least four expressive extensions and their parsing implications.\\
16 & Error Handling & Ch. 15 & Define diagnostic quality metrics and regression tests for three syntactic fault classes.\\
17 & Practical Parser Writing and Usage & Ch. 16 & Produce a parser-engineering specification with APIs, error contracts, AST ownership, tests and performance targets\\
18 & Annotated Bibliography & Ch. 17 & Prepare an annotated reading register: contribution, model, assumptions, complexity, implementation evidence.\\
\bottomrule\end{longtable}\normalsize
\section{Prerequisite graph}
Foundational track: 1 $\to$ 2 $\to$ 3; branch 3 $\to$ 5 (automata) and 3 $\to$ 4 (CYK); 5 $\to$ 6 $\to$ 8, and 4+6 $\to$ 7; 7+8 $\to$ 9. Advanced paths: 7+9 $\to$ 10+11; 5+7 $\to$ 12+13; 4+11 $\to$ 14; 2+3 $\to$ 15; 6+7+8+9 $\to$ 16; 8+9+16 $\to$ 17; 18 is a continuous research literacy strand, not an algorithmic prerequisite.\par
\section{Competency traceability}
Each study guide has an identified chapter and evidence tasks. Core practical examinations: grammar/automata workbook (Ch. 2/5), CYK/Earley recognizers (Ch. 4/7), LL(1)/LR conflict diagnostics (Ch. 8/9), parsing security review (Ch. 16/17), and a final end-to-end compiler (Ch. 17).\par
\PTSource{Contents and corresponding chapter topics, 2nd edition (2008).}

\end{document}
